% !TeX root = ../colloquium-2.tex

\phantomsection\label{sec_Proofs01}
\begin{proof}[Доказательство теоремы о нуле непрерывной на отрезке функции. Доказательство теоремы Коши о непрерывной функции.]
    \subsection*{Доказательство теоремы о нуле непрерывной на отрезке функции.}
    \begin{note}[amber]
        Пусть функция $f \in C \ ([a, b])$, и $f(a) \cdot f(b) < 0$. Тогда $ \exists \ c \in (a, b)$, такой что $f(c) = 0$. \\
        \begin{note}[purple]
            Доказательство.По условию, если $f(a) < 0$, то $f(b) > 0$, принимает значение разных знаков в точках $a$ и $b$. \\
            Разобьем отрезок $[a, b]$ пополам, если $f(\frac{a + b}{2}) = 0$, то все доказано, если нет,то пусть, например, $f(\frac{a + b}{2}) < 0$, обозначим $\frac{a + b}{2} = a_{1}$, $b = b_{1}$, для отрезка $[a_{1}, b_{1}]$ проделаем ту же операцию, \\
            возьмем точку $(\frac{a_{1} + b_{1}}{2})$, найдем $f(\frac{a_{1} + b_{1}}{2})$, то все доказано, если это значение не равно нулю, то из отрезков $[a_{1}, \frac{a_{1} + b_{1}}{2}]$ и $[\frac{a_{1} + b_{1}}{2}, b_{1}]$ выберем тот, на концах которого функция $f$ принимает значения разных знаков. \\
            Обозначим через $a_{2}$ левый конец этого отрезка, а через $b_{2}$ - правый азобьём этот отрезок
            пополам и т.д.. Если ни в какой из середин отрезков не получится нулевое значение функции, тот этот процесс будем продолжать до бесконечности и получим последовательность
            стягивающихся отрезков. Пересечение этих отрезков состоит, как мы знаем, из одной точки, которую обозначим через c. Последовательность левых концов построенных отрезков $\{a_n\}^{+\infty}_{n = 1}$ и последовательность правых концов $\{b_n\}^{+\infty}_{n = 1}$ стремятся к c,
            поэтому, в силу непрерывности функции на отрезке $[a, b] \ f (a_n) \rightarrow f(c), f(b_n) \rightarrow f(c)$ и при этом $f(a_n) \cdot f(b_n) < 0$ по построению. Используя предельный переход в неравенствах и арифметику предела,
            имеем: $f^2(c) = \displaystyle{}\lim_{n \to \infty} f(a_n) \cdot f(b_n) \leq 0$, поэтому $f(c) = 0$, чем все и доказано. 
        \end{note}
    \end{note}

    \subsection*{Доказательство теоремы Коши о непрерывной функции.}
    \begin{note}[amber]
        Пусть функция $f$ определена и непрерывна на отрезке $[a,b], m = \displaystyle{}\inf_{a \leq x\leq b} f(x), M = \displaystyle{}\sup_{a\leq x \leq b} f(x).$ Тогда для любого
        числа C $\in [m,M]$ существует такая точка c $\in [a,b]$, что $f(c) = C.$\\
        \begin{note}[purple]
           Докозательство. Если $C = m$ или $C = M$, то всё доказано. Если С $\in (m,M)$, то функция $g:= f - C$ непрерывна на отрезке $[x_1,x_2]$(отрезок неорентированный) и $g(x_1) = m - C < 0$, а $g(x_2) = M - C > 0$, то есть для функции $g$ на отрезке $[x_1,x_2]$ выполнены все условия теоремы  о нуле непрерывной на отрезке функции
           . Тогда $\exists c \in [x_1,x_2] : g(c) = 0$, то есть $f(c) = C$.
        \end{note}
    \end{note}
\end{proof}
